% Двете конфигурации за измерване, поставени една над друга.
%
% Рисунката трябва да носи едно твърдение преди да е прочетен ред текст: втората
% конфигурация допуска по-тесен кръг от изводи. Затова двата панела завършват с еднакво
% широка лента „обхват“. При основната конфигурация цялата лента е запълнена, при
% вторичната запълнени са само около две пети отляво, а останалото е празен пунктирен
% правоъгълник.
% Сравнението е геометрично и не зависи от четенето.
%
% Черно-бяла рисунка. Разграничението е по форма, по дебелина на линията и по сив тон:
% активните единици са gray!25, поясненията gray!10, ограниченията са в рамка с удебелен
% контур, а липсващата възможност е зачертана с кръст. Няма нито един цвят.
\begin{figure}[htbp]
\centering
\begin{tikzpicture}[
    % Вътрешните координати стигат до около 15,3 по хоризонтала. Мащабът 0,9 ги свива до
    % под 13,8 cm, тоест под ширината на текста 14,5 cm.
    scale=0.9, transform shape,
    % preamble.tex дефинира note два пъти: веднъж като чист текст и по-късно като
    % пунктирана сива кутия. Втората дефиниция е тази, която важи в дисертацията, тоест
    % без този ред всеки надпис тук се огражда с кутия, рамките се застъпват със
    % стрелките и с пунктирните пространства, а разликата между пояснение и ограничение
    % изчезва. Тук note се връща към чист текст, но само за тази рисунка.
    note/.style={align=left, font=\fontsize{8}{10}\selectfont},
    unit/.style={proc, minimum width=4.6cm, minimum height=0.95cm, fill=gray!25,
                 font=\fontsize{8}{9.5}\selectfont},
    small unit/.style={unit, minimum width=3.6cm},
    addr/.style={box, minimum width=4.6cm, minimum height=0.85cm, fill=gray!10,
                 font=\fontsize{7.5}{9}\selectfont},
    nslab/.style={note, align=center, font=\fontsize{8}{9.5}\selectfont\bfseries},
    prop/.style={box, text width=4.2cm, align=center, fill=gray!10, minimum height=1.05cm,
                 font=\fontsize{7.5}{9}\selectfont},
    lim/.style={box, very thick, text width=4.0cm, align=left, fill=gray!10,
                minimum height=1.70cm, font=\fontsize{7.5}{9}\selectfont},
    scopeok/.style={box, anchor=west, minimum height=1.25cm, fill=gray!25, align=center,
                    font=\fontsize{7.5}{9}\selectfont},
    scopeno/.style={box, anchor=west, minimum height=1.25cm, densely dashed, align=center,
                    font=\fontsize{7.5}{9}\selectfont},
    head/.style={align=center, font=\fontsize{9.5}{11.5}\selectfont\bfseries},
    sub/.style={note, align=center, text width=13.2cm, font=\fontsize{7.5}{9}\selectfont},
    tag/.style={note, align=center, font=\fontsize{7}{8.5}\selectfont, inner sep=1.5pt},
    panel/.style={draw, thick, rounded corners=3pt, inner sep=8pt},
  ]

  % =============================================================== основна конфигурация
  \node[head] (mt) at (7.70, 1.05) {Основна конфигурация (Linux): двойка мрежови
                                    пространства};
  \node[sub]  (ms) at (7.70, 0.45) {Две мрежови пространства на една машина, свързани с
                                    двойка \code{veth}.};

  % Пространството на генератора носи два адреса. Това не е украса: без втори адрес
  % миграцията на връзка няма как да се случи, тоест твърдението за нея е непроверимо.
  \node[nslab] (nsg) at (2.90, -0.30) {мрежово пространство \code{gen}};
  \node[unit]  (gen) at (2.90, -1.20) {генератор\\ h2load, wrk2, k6};
  \node[addr]  (adg) at (2.90, -2.50) {\code{veth-gen}: 10.0.0.2\\
                                       втори адрес: 10.0.0.3};
  \node[grp, fit=(nsg)(gen)(adg)] (gg) {};

  \node[nslab] (nss) at (12.50, -0.30) {мрежово пространство \code{srv}};
  \node[unit]  (srv) at (12.50, -1.20) {сървър\\ \code{coroute}};
  \node[addr]  (ads) at (12.50, -2.50) {\code{veth-srv}: 10.0.0.1};
  \node[grp, fit=(nss)(srv)(ads)] (gs) {};

  % Двете посоки се рисуват като две отделни стрелки, защото \code{netem} се задава на
  % изхода на всяко \code{veth} поотделно. Една двупосочна стрелка би подсказала точно
  % грешката, която разделът за loopback описва.
  \draw[flow] (gg.east |- gen) -- (gs.west |- gen);
  \node[tag] at (7.70, -0.68) {\code{netem} към сървъра\\
                               закъснение, загуба, \code{limit}};

  \node[box, fill=gray!10, font=\fontsize{7}{8.5}\selectfont] (vp) at (7.70, -1.88)
        {двойка \code{veth}};

  \draw[flow] (gs.west |- ads) -- (gg.east |- ads);
  \node[tag] at (7.70, -3.02) {\code{netem} към клиента\\ зададено поотделно};

  % Трите свойства, заради които тази конфигурация носи централното твърдение.
  \node[prop] (p1) at (2.90, -4.30) {истински MTU 1500 и истинска опашка \code{qdisc}};
  \node[prop] (p2) at (7.70, -4.30) {истински IP път, а не \code{lo}};
  \node[prop] (p3) at (12.50, -4.30) {два адреса на клиента, миграцията е проверима};

  \node[note, anchor=north west, align=left, text width=14.1cm,
        font=\fontsize{7.5}{9.5}\selectfont] (nt) at (0.62, -5.15)
       {Ядрата са изолирани, генераторът и сървърът не делят ядро.\\
        \code{netem} стои на изхода на всяко \code{veth} поотделно, тоест зададеното
        закъснение е закъснение за посока, а не за обръщане.\\
        Дължината на опашката \code{limit} се задава изрично, защото подразбиращите се
        1000 пакета превръщат зададената загуба в препълване.};

  \node[scopeok, text width=13.6cm, minimum width=14.15cm] (bar1) at (0.62, -7.60)
       {обхват: сравнение между HTTP/1.1, HTTP/2 и HTTP/3, миграция на QUIC, влошена
        мрежа, горни персентили};

  % Двете невидими точки държат двата панела с еднаква ширина. Без тях рамките се
  % получават с няколко милиметра разлика, което се чете като небрежност.
  \node[panel, fit={(mt)(ms)(gg)(gs)(p1)(p2)(p3)(nt)(bar1)(0.30,-1.20)(15.20,-1.20)}]
        (mpanel) {};

  % ============================================================= вторична конфигурация
  \node[head] (st) at (7.70, -9.20) {Вторична конфигурация (Windows): loopback};
  \node[sub]  (ss) at (7.70, -9.75) {Генераторът и сървърът са на един хост и говорят
                                     през \code{lo}.};

  % Кутиите тук са по-тесни от горните, за да остане място вдясно за това, което липсва.
  \node[nslab] (hl) at (6.10, -10.40) {един хост};
  \node[small unit] (gen2) at (2.40, -11.20) {генератор\\ h2load, k6};
  \node[box, fill=gray!10, font=\fontsize{8}{9.5}\selectfont] (lo) at (6.10, -11.20)
       {\code{lo}\\ 127.0.0.1};
  \node[small unit] (srv2) at (9.80, -11.20) {сървър\\ \code{coroute}};
  \draw[flow, {Stealth}-{Stealth}] (gen2) -- (lo);
  \draw[flow, {Stealth}-{Stealth}] (lo) -- (srv2);
  \node[grp, fit=(hl)(gen2)(lo)(srv2)] (gh) {};

  % Зачертаната кутия е единственият елемент с кръст в рисунката. Тя казва какво липсва,
  % и затова стои извън пунктираната рамка на хоста. Вътре в нея стои една дума, защото
  % кръстът минава през текста и по-дълъг надпис не би се четял.
  \node[box, densely dashed, minimum width=2.4cm, minimum height=0.9cm] (sh)
        at (13.40, -11.20) {};
  \draw[very thick] (sh.north west) -- (sh.south east);
  \draw[very thick] (sh.south west) -- (sh.north east);
  % Надписът се поставя след кръста и с бял фон, иначе линиите минават през буквите.
  \node[fill=white, inner sep=1.5pt, font=\fontsize{8}{9.5}\selectfont]
        at (sh) {\code{netem}};
  \node[tag, below=3pt of sh] (shl) {няма програмируемо\\
                                     оформяне на трафика,\\
                                     няма \code{tc}};

  % Ограниченията се назовават тук, а не се оставят за раздела с изводите.
  \node[lim] (l1) at (2.85, -13.75) {\textbf{MTU 65536.} TCP получава сегменти от
        64\,KB, а QUIC остава при 1200 до 1450 байта. Сравнението наказва HTTP/3.};
  \node[lim] (l2) at (7.70, -13.75) {\textbf{Обръщане в микросекунди.} Скриването на
        закъснение от мултиплексирането става невидимо.};
  \node[lim] (l3) at (12.55, -13.75) {\textbf{Един адрес.} Миграцията на връзка е
        непроверима по определение.};

  % Лентата има същата обща ширина като горната. Запълнената част е обхватът, празната е
  % това, което тази конфигурация не може да поддържа.
  \node[scopeok, text width=5.0cm, minimum width=5.60cm] (bar2) at (0.62, -15.55)
       {обхват: микроизмервания на маршрутизатора, цена на сериализацията, преносимост};
  \node[scopeno, text width=8.0cm, minimum width=8.55cm] (bar3) at (6.22, -15.55)
       {извън обхвата: сравнение между протоколите, миграция на връзка, влошена мрежа,
        горни персентили};

  \node[panel, fit={(st)(ss)(gh)(sh)(shl)(l1)(l2)(l3)(bar2)(bar3)
                    (0.30,-11.20)(15.20,-11.20)}] (spanel) {};

\end{tikzpicture}
\caption{Двете конфигурации за измерване. Горе: основната конфигурация под Linux, две
мрежови пространства на една машина, свързани с двойка \code{veth}, с \code{netem} на
изхода на всяка посока поотделно и с изрично зададена дължина на опашката. Тя дава
истински MTU, истинска опашка, истински IP път и два адреса на клиента, тоест миграцията
на QUIC е проверима. Долу: вторичната конфигурация под Windows, генератор и сървър на един
хост през \code{lo}, без програмируемо оформяне на трафика. Трите ограничения са
отбелязани изрично. Долните ленти имат еднаква обща ширина и сравняват обхвата на
твърденията, които всяка конфигурация допуска.}
\label{fig:measurement_topology}
\end{figure}
